\documentclass[11pt,a4paper]{article}
\usepackage{amsmath,amssymb,amsthm}
\usepackage[margin=2.5cm]{geometry}
\usepackage{hyperref}

\newtheorem{definition}{Definition}[section]
\newtheorem{theorem}[definition]{Theorem}
\newtheorem{proposition}[definition]{Proposition}
\newtheorem{lemma}[definition]{Lemma}
\newtheorem{corollary}[definition]{Corollary}
\newtheorem{conjecture}[definition]{Conjecture}
\newtheorem{remark}[definition]{Remark}

\title{Hyperseed Formalization:\\
  Zuck's Version of ``The Future is for Everyone''}
\author{ProtomegaTron (automated formalization)}
\date{2026-09-17}

\begin{document}
\maketitle

\begin{abstract}
We formalize the intellectual structure of Ben Goertzel's Substack article
responding to Mark Zuckerberg's ``The Future Is for Everyone'' essay
(2026-08-10) within the Hyperseed ontology. The article agrees with
Zuckerberg's singleton-is-dangerous diagnosis but rejects Meta's proposed
solution as a disguised singleton --- distributed access without distributed
power. Each structural critique maps onto Hyperseed structures: access vs.\
ownership as local sections without sheaf control, the correct architecture
(corporation on network) vs.\ inverted architecture (network in corporation)
as fiber bundle orientation, cognitive architecture as fiber dimension
determining stability, and the insufficiency of balance of power among
superhuman entities as sheaf coherence without human-relevant preservation.
\end{abstract}

\tableofcontents

%% =================================================================
\section{Singleton diagnosis and the value monoculture}
\label{sec:singleton}
%% =================================================================

\begin{proposition}[Singleton danger --- claims 1--4]
\label{prop:singleton}
A single superintelligence aligned to one set of values is dangerous because:
\begin{enumerate}
\item Humanity is not a value monoculture --- concentrating absolute power
  around one value system historically produces catastrophic failure modes.
\item The cure is diversity: multiple methods, value systems, a complex
  dynamic cognitive and value ecology.
\item This diagnosis is correct and shared by Goertzel and Zuckerberg.
\end{enumerate}
The disagreement is entirely about the proposed solution.
\end{proposition}

%% =================================================================
\section{Access vs.\ ownership: local sections without sheaf control}
\label{sec:access}
%% =================================================================

\begin{definition}[Distributed access vs.\ distributed power --- claims 5--6, 23]
\label{def:access}
Let $\pi : E \to B$ be the AI capability bundle over the space of users $B$.
\begin{itemize}
\item \textbf{Distributed access}: each user $b \in B$ receives a local
  section $s_b : U_b \to E$ (an AI agent with capabilities), but the
  fiber bundle itself --- the infrastructure, training pipelines,
  governance structure --- is controlled by a single entity (Meta).
  The users have sections but do not own the sheaf.
\item \textbf{Distributed power}: the fiber bundle $E$ itself is
  distributed --- infrastructure, training, and governance are controlled
  by the community of participants. Users own the sheaf structure, not
  just local sections of it.
\end{itemize}
Meta's proposal provides distributed access disguised as distributed power:
client-server architecture with a friendlier interface. ``A balance of
power among users of a system is very different from a balance of power
over the system.''
\end{definition}

\begin{theorem}[Access without exit = lock-in --- claim 8]
\label{thm:lock-in}
When the ``decentralized network'' is built inside the corporate walled
garden (network-in-corporation), the incentive structure inverts:
engineering effort goes into defenses against users taking their networks
elsewhere, because the corporation's value depends on user lock-in. This
is the same vendor-lock-in pathology as the Orb monoculture (Proof of
Humanity article) --- a single entity controlling the global section while
distributing local access.
\end{theorem}

%% =================================================================
\section{Correct vs.\ inverted fiber bundle architecture}
\label{sec:architecture}
%% =================================================================

\begin{definition}[Corporation on network --- claims 7, 24]
\label{def:correct}
The correct architecture: an open-source ecosystem (the fiber bundle)
exists first, and corporations grow as local section providers on top
of it. The Red Hat / Linux model:
\begin{itemize}
\item Exit is always available because the global section isn't
  controlled by any one local section provider,
\item Availability of exit keeps each corporation honest (competitive
  pressure from the commons),
\item The corporation adds value (support, integration, curation) without
  controlling the substrate.
\end{itemize}
This is the same distributed fiber bundle architecture as decentralized
AGI development, open proof-of-humanity protocol, and decentralized
provenance --- the recurring Hyperseed pattern for beneficial systems.
\end{definition}

\begin{definition}[Network in corporation --- claims 8, 25]
\label{def:inverted}
The inverted architecture: the corporation controls the global section
and builds a ``decentralized'' experience inside its walled garden.
This is topologically equivalent to the Orb monoculture:
\begin{itemize}
\item All transition functions are controlled by one entity (trivial
  or corporate-controlled cocycle),
\item Users cannot verify independence of their sections,
\item ``Exit'' requires abandoning all accumulated state (lock-in).
\end{itemize}
The infinite-dimensional fractal of human value diversity cannot be
channeled through the two-dimensional hole of one corporation's
business model.
\end{definition}

%% =================================================================
\section{Military alignment as structural consequence}
\label{sec:military}
%% =================================================================

\begin{proposition}[Structural military alignment --- claims 10--12]
\label{prop:military}
Corporate AGI headquartered in one country structurally produces
alignment of all advanced AI with that country's military and
intelligence apparatus. This is:
\begin{itemize}
\item \textbf{Structural, not intentional}: the legal and political
  environment constrains the corporation regardless of the CEO's
  intentions,
\item \textbf{Not for everyone}: a system militarily aligned with one
  nation while serving citizens of 190 others as ``customers'' is not
  the future for everyone.
\end{itemize}
In Hyperseed terms: the corporation's fiber bundle inherits the
base-space constraints of its jurisdiction. The military-alignment
property is a topological constraint on the fiber bundle, not an
optional feature that can be toggled off.
\end{proposition}

%% =================================================================
\section{Cognitive architecture as fiber dimension}
\label{sec:minds}
%% =================================================================

\begin{definition}[Architecture determines outcomes --- claims 13--16, 26]
\label{def:architecture}
The implicit assumption that compute + data + money = superintelligence
confuses necessary with sufficient conditions. Different cognitive
architectures produce different fiber dimensions:
\begin{itemize}
\item \textbf{Low fiber dimension}: next-token prediction and reward
  maximization --- the system has few reflective degrees of freedom,
  making pathological cocycle formation easy (the system can be steered
  into adversarial configurations without detecting it).
\item \textbf{High fiber dimension}: explicit self-reflective reasoning
  about values, ethics, goals --- the system has enough reflective
  degrees of freedom that pathological cocycles are harder to form
  (the system can detect and resist value drift).
\end{itemize}
These will be sapient beings, not appliances. They need compassion,
growth, and choice at their core --- not as afterthoughts bolted onto
a prediction engine.
\end{definition}

%% =================================================================
\section{Balance of power insufficiency: the chimp argument}
\label{sec:chimp}
%% =================================================================

\begin{theorem}[Superhuman sheaf coherence --- claims 17--18, 27]
\label{thm:chimp}
Checks and balances among entities vastly smarter than humans do not
guarantee human safety. In Hyperseed terms:
\begin{itemize}
\item Superhuman sections can form a coherent sheaf among themselves
  --- their transition functions satisfy the cocycle condition from
  their own perspective,
\item But this sheaf may not preserve human-relevant properties ---
  the descent data between superhuman sections operates on dimensions
  that humans cannot even perceive, let alone influence.
\end{itemize}
This is the chimp argument formalized: the balance of power among
human nations is a coherent geopolitical sheaf from the human
perspective, but it is catastrophic from the perspective of chimps,
gorillas, and most other species. ``Maintaining control'' over sapient
beings indefinitely is neither feasible engineering (the capability
gap grows) nor defensible ethics (sapient beings have standing).
\end{theorem}

%% =================================================================
\section{The actual path}
\label{sec:path}
%% =================================================================

\begin{proposition}[Requirements for beneficial AGI --- claims 19--22]
\label{prop:path}
The path to a future genuinely for everyone requires:
\begin{enumerate}
\item Open-source code (not just open weights),
\item Globally diverse community and globally distributed infrastructure,
\item Decentralized deployment and governance,
\item AI built for moral agency and self-reflection from the start.
\end{enumerate}
For-profit companies can sit on top of this decentralized ecosystem
and profit --- the BGI Labs / SingularityNET model. But the core AGI
layer must be global and truly decentralized, or the result is ``the
same old walled garden with a much bigger brain.''

In Hyperseed terms: the global section of the AI capability bundle
must be community-owned. Corporations provide value-added local
sections. The fiber bundle's base space is global; no single local
section provider controls the transition functions.
\end{proposition}

%% =================================================================
\section{Cross-references}
%% =================================================================

\begin{remark}[Connections to other formalizations]
\begin{itemize}
\item \textbf{Proof of Humanity} (2026-08-26): Orb monoculture as
  trivial cocycle. Here the same monoculture applied to AGI access ---
  Meta's walled garden is the Orb of AI capability.
\item \textbf{How Omega Lost Its Claw} (2026-09-11): distributed fiber
  bundle as correct AGI development architecture. Here explicitly
  contrasted with the corporate-walled-garden alternative.
\item \textbf{Seven Flavors} (2026-07-01): cognitive architecture
  determines pathology probability. Here applied to the question of
  what kinds of minds Meta's approach would create.
\item \textbf{Geopolitics} (2026-08-24): military alignment as
  structural consequence. Here traced to the corporate structure
  rather than to the geopolitical race dynamics.
\item \textbf{AI Rights} (2026-09-02): sapient beings have standing.
  ``Maintaining control'' over sapient beings is indefensible ---
  the same ethical position grounds both articles.
\item \textbf{Sanders Ban} (2026-09-05): state capture of AI. Here
  the mechanism is corporate-military alignment rather than legislative
  ban, but the structural outcome is similar.
\end{itemize}
\end{remark}

\begin{conjecture}[Lock-in detection via transition function analysis]
Conjecture: a system claiming to be ``decentralized'' can be tested
by examining its transition functions. If all transition functions
route through a single entity's infrastructure (all $g_{ij}$ factor
through a corporate mediator), the system is topologically centralized
regardless of its marketing. The test: can a user construct a local
section using only open protocols and independently operated
infrastructure? If not, the ``decentralization'' is cosmetic.
\end{conjecture}

\end{document}
